% Contexto e instrucciones literales de las páginas 1 y 2 de EXA-0071.
% Fuente: Selectividad/sel_2009_matematicas_aplicadas/mat apli CC SS05.pdf
% SHA-256: ca51ab435c7584a8e047e19d0826f3273f3e88aed6110ccd00e46a40e85518b1
% No se reproduce el logotipo ni el marco de la cabecera.
\section*{Contexto del examen original}
\textbf{UNIVERSIDADES DE ANDALUCÍA}\par
\textbf{PRUEBA DE ACCESO A LA UNIVERSIDAD}\par
CURSO 2008-2009\par
\textbf{MATEMÁTICAS APLICADAS A LAS CIENCIAS SOCIALES II}

\textbf{Instrucciones:}
\begin{enumerate}[label=\textbf{\alph*)}]
  \item Duración: 1 hora y 30 minutos.
  \item Elija una de las dos opciones propuestas y conteste los ejercicios de la opción elegida.
  \item En cada ejercicio, parte o apartado se indica la puntuación máxima que le corresponde.
  \item Se permitirá el uso de calculadoras que no sean programables, gráficas ni con capacidad para almacenar o transmitir datos.
  \item Si obtiene resultados directamente con la calculadora, explique con detalle los pasos necesarios para su obtención sin su ayuda. Justifique las respuestas.
\end{enumerate}
\noindent\textbf{Nota de la colección:} Se reúnen ambas opciones para consulta. La regla anterior pertenece al examen original. Curso 2008-2009 impreso;convocatoria,tipo,reserva y modelo no indicados. El05del nombre no acredita modelo. PDF escaneado: páginas1y2 de examen y página3 de criterios específicos. Sin tabla normal asociada. Cada opción contiene tres ejercicios:3,3y4puntos;el tercero conserva Parte I y Parte II de2puntos cada una,con numeración original y clasificación en probabilidad e inferencia.
\par\medskip
